
\section{評価結果}
\subsection{RQ1の結果}
手法ごとの修正率とプロンプト長を\textbf{表\ref{tab:main}}に示す．
提案手法の修正率は68.6\%であった．リンクのみを与えるBL0が0.0\%，同じリンクから本文を取得できるBL-Webが45.7\%であることから，分岐点はリンクを与えることではなく仕様本文に到達できることにある．しかしBL-Webは本文に到達してなお提案手法に22.9ポイント及ばない．先行研究の構成であるDsLsは31.4\%であり，提案手法はその2.2倍の修正率を約8分の1のプロンプト長で達成した．

\begin{mdframed}[linewidth=1pt,linecolor=black]{RQ1への回答：}
仕様本文に到達できるBL-Webに対して22.9ポイント，先行研究の構成に対して37.2ポイント修正率を向上させた．プロンプト長は先行研究の約8分の1である．
\end{mdframed}

\subsection{RQ2の結果}
アブレーションの結果，除去による低下が最も大きかったのは対応付け先からの展開であり，68.6\%から14.3\%へ54.3ポイント低下した．差分カードへの整形の除去は11.5ポイント，リクエスト事実への分割の除去は5.7ポイントの低下であった．

A4は差分ノードそのものは提案手法と同一のまま保持している．したがって，旧仕様と新仕様の対応関係を与えるだけでは14.3\%にとどまり，そこから新仕様で必要となる要素を構造的に引き寄せることが修正率を押し上げていると言える．

\begin{mdframed}[linewidth=1pt,linecolor=black]{RQ2への回答：}
対応付け先からの展開が支配的であり，これを除去すると54.3ポイント低下する．差分カードへの整形は11.5ポイント，事実への分割は5.7ポイントの寄与である．
\end{mdframed}

\subsection{同一の修正を要する誤用群における比較}
対象データセットには，\texttt{/v1.0/}を\texttt{/v1.1/}へ更新した上で\texttt{t}，\texttt{sign}，\texttt{nonce}の各ヘッダをHMAC-SHA256による署名とともに付与するという，\textbf{同一の修正を要する}誤用群が含まれる．この誤用群における条件ごとの成功率を\textbf{表\ref{tab:class}}に示す．

提案手法が93.3\%であるのに対し，本文に到達できるBL-Webは40.0\%であった．BL-Webはこの誤用群の一部では提案手法を上回る一方で他では全く修正できておらず，同一の修正を要するにもかかわらず結果が安定しない．A4は0.0\%であり，展開なしには一件も修正できていない．

\begin{table}[t]
\centering
\caption{同一の修正を要する誤用群における成功率}
\label{tab:class}
\scriptsize
\begin{tabular}{@{}lr@{}}
\toprule
手法 & 成功率 (\%) \\
\midrule
提案手法 & 93.3 \\
A3 & 86.7 \\
A2 & 66.7 \\
BL-Web & 40.0 \\
DsLs & 13.3 \\
A4 & 0.0 \\
BL0 & 0.0 \\
\bottomrule
\end{tabular}
\end{table}
